% ================================================================
% Section: Deployment on Resource-Constrained Environments
% From main (2).tex — Section VI
% ================================================================

\section{Deployment on Resource-Constrained Environments}
\label{sec:deployment}

% ----------------------------------------------------------------
\subsection{On-Device and Mobile Deployment}
\label{subsec:deployment_mobile}

On-device deployment eliminates the latency, cost, and privacy exposure of cloud API inference. It enables AI functionality in offline and low-bandwidth environments and is a prerequisite for truly private personal assistants. The primary challenges are memory (RAM and storage), sustained inference throughput under thermal constraints, and battery impact.

\subsubsection{Apple Intelligence (AFM)}

Apple's Foundation Models \cite{apple2024intelligence,gunter2025afm} are the most fully documented on-device \gls{slm} deployment. The 2025 report \cite{gunter2025afm} describes a $\sim$3B on-device model, optimized for Apple silicon through KV-cache sharing (dividing the model so the second block reuses the first block's cache, a reported 37.5\% memory saving) and 2-bit quantization-aware training, paired with a Parallel-Track Mixture-of-Experts server model on Private Cloud Compute. The on-device model is reported competitive with Qwen3-4B and Gemma-3-4B on its target tasks; the server model with Llama~4 Scout. The third-generation framework \cite{apple2026afm3} adds a natively multimodal, sparsely activated on-device model and opens the same API surface to third-party cloud models, so that an application calls the local \gls{slm} by default and escalates to a frontier \gls{llm} only on demand. Google's Gemini Nano and Samsung's Gauss \cite{samsung2023gauss} occupy the same design point on Android and Galaxy devices.

\subsubsection{Samsung Gauss}

Samsung Gauss \cite{samsung2023gauss} is a family of language models for Galaxy devices, including a ${\sim}1.7$B on-device model targeting the Gauss1 chip and Exynos 2400 \gls{npu}. Samsung reports that Gauss enables on-device email summarization, translation, and writing assistance without requiring internet connectivity.

\subsubsection{General Mobile Deployment Considerations}

\glspl{slm} targeting mobile deployment must navigate three constraints. (1) \emph{Memory}: Consumer smartphones typically provide 6--12~GB RAM, of which the OS uses 2--4~GB, leaving 4--8~GB for applications. A 3B model at Q4 requires ${\sim}2$~GB, comfortable for most flagship devices. (2)~\emph{Thermal}: Sustained inference on mobile \glspl{npu} is throttled by thermals; realistic sustained throughput is lower than peak throughput benchmarks suggest. (3)~\emph{Battery}: Gemma~3 on Android reportedly consumes ${\sim}0.75\%$ battery per prompt on Pixel~9 Pro, enabling hundreds of interactions per charge. The fundamental constraint in edge deployment is memory bandwidth rather than compute: mobile devices offer 50--90~GB/s of memory bandwidth versus 2--3~TB/s for datacenter GPUs \cite{xu2024device}, making INT4 quantization not merely a storage optimization but a critical performance enabler.

% ----------------------------------------------------------------
\subsection{Edge and IoT Deployment}
\label{subsec:deployment_edge}

Below smartphone-class hardware lies a diverse ecosystem of edge devices: Raspberry Pi single-board computers, NVIDIA Jetson modules, and microcontrollers. \glspl{slm} targeting this tier typically require:

\begin{itemize}
    \item \textbf{Sub-1B or extreme quantization:} A Raspberry Pi~5 (8~GB RAM, Cortex-A76 CPU) can run Q4 quantized 1B models at 10--20 tokens/second. SmolLM2-135M at Q8 runs at $>$30~tok/s on Raspberry Pi~5. Benchmarking shows that llamafile achieves up to $4\times$ higher throughput and 30--40\% lower power consumption than Ollama for equivalent models on single-board computers \cite{costa2025sbc}.

    \item \textbf{Jetson Orin:} NVIDIA's Jetson Orin NX (16~GB) with TensorRT-LLM can serve 7B quantized models at ${\sim}30$~tok/s, enabling real-time local inference for industrial and robotics applications.

    \item \textbf{Microcontrollers (MCU):} Sub-100M models using aggressive quantization (INT4/INT2) can run on ARM Cortex-M series MCUs with 512~KB--2~MB SRAM for specialized tasks like keyword detection and sensor data analysis.
\end{itemize}

% ----------------------------------------------------------------
\subsection{Deployment Frameworks}
\label{subsec:deployment_frameworks}

A mature ecosystem of open-source frameworks enables \gls{slm} deployment across hardware targets. Table~\ref{tab:deployment_frameworks} provides a structured comparison.

\begin{table}[ht]
\caption{Key SLM deployment frameworks, their primary targets, and characteristics.}
\label{tab:deployment_frameworks}
\centering
\footnotesize\adjustbox{max width=\linewidth}{%
\begin{tabular}{@{}lllp{5.5cm}@{}}
\toprule
\textbf{Framework} & \textbf{Primary Target} & \textbf{Quantization} & \textbf{Notes} \\
\midrule
llama.cpp \cite{gerganov2023llamacpp}   & CPU, Metal, CUDA, ROCm & GGUF (Q2--Q8)   & De facto standard; minimal dependencies; ecosystem includes Ollama, llamafile \\
ONNX Runtime \cite{microsoft2019onnx} & GPU, CPU, NPU, WASM & INT4/INT8       & Cross-platform; DirectML/CoreML/CUDA backends; Phi-3-mini at 35~tok/s on Snapdragon NPU \\
MLX \cite{mlx2023}             & Apple Silicon          & 4-bit           & Unified memory; eliminates GPU transfers; 55+ tok/s for LLaMA-3.2-3B on M3~Pro \\
TensorRT-LLM \cite{nvidia_tensorrt2024} & NVIDIA GPUs   & INT4/INT8/FP16  & Continuous batching; kernel fusion; 2--3$\times$ over PyTorch for 7B models \\
ExecuTorch \cite{meta2024executorch} & Mobile (iOS, Android) & INT4          & 50~KB base runtime; 12+ hardware backends; deployed in Instagram, WhatsApp \\
MLC-LLM \cite{mlc2024}        & iOS, Android, web, MCU & 4-bit          & TVM compilation; WebGPU support; universal platform coverage \\
WebLLM \cite{webllm2024}       & Web browsers           & 4-bit WASM     & WebGPU acceleration; 80\% of native performance; OpenAI-compatible API \\
\bottomrule
\end{tabular}%
}
\end{table}

% ----------------------------------------------------------------
\subsection{Latency Comparison: SLMs vs.\ LLMs}
\label{subsec:deployment_latency}

Table~\ref{tab:latency_compare} illustrates the practical latency implications of \gls{slm} vs.\ \gls{llm} deployment for three representative real-world scenarios.

\begin{table}[ht]
\small
\caption{Latency and cost comparison: SLMs vs.\ LLMs for representative use cases.}
\label{tab:latency_compare}
\centering
\footnotesize\adjustbox{max width=\linewidth}{%
\begin{tabular}{@{}lp{2.5cm}p{2.5cm}p{2.5cm}@{}}
\toprule
\textbf{Metric} & \textbf{Real-time Chatbot} & \textbf{Code Completion} & \textbf{Edge IoT} \\
\midrule
\multicolumn{4}{@{}l}{\textit{SLM (3B, Q4, consumer GPU)}} \\
TTFT               & 80--150 ms    & 30--80 ms     & 50--120 ms  \\
Throughput         & 200--400 tok/s & 300--500 tok/s & 100--300 tok/s \\
Memory             & ${\sim}2$~GB  & ${\sim}2$~GB  & ${\sim}1$~GB \\
Cost/1M tokens     & \$0.05--0.2   & \$0.05--0.2   & \$0 (on-device) \\
\midrule
\multicolumn{4}{@{}l}{\textit{LLM (70B, FP16, multi-GPU)}} \\
TTFT               & 500--2000 ms  & 200--800 ms   & Not feasible \\
Throughput         & 30--80 tok/s  & 60--150 tok/s & N/A \\
Memory             & ${\sim}140$~GB & ${\sim}140$~GB & N/A \\
Cost/1M tokens     & \$3--10       & \$3--10        & N/A \\
\midrule
\multicolumn{4}{@{}l}{\textit{Frontier API (GPT-4o, cloud)}} \\
TTFT               & 300--1500 ms  & 200--600 ms   & Requires internet \\
Throughput         & 80--200 tok/s & 80--200 tok/s & N/A \\
Cost/1M tokens     & \$5--15       & \$5--15        & N/A \\
\bottomrule
\end{tabular}%
}
\end{table}

The data in Table~\ref{tab:latency_compare} illustrate a 5--20$\times$ cost advantage and comparable or superior latency for \glspl{slm} in chatbot and code-completion scenarios, with an absolute deployment advantage in offline edge contexts where \glspl{llm} are simply infeasible.
